% Generated by tools/build_v5.py; edit the assembly rules, not this file.
\section{Asymmetric screening above exponent two}\label{app:above-two}
The upper proof is uniform over the stated envelope and global rarity
class. The matching lower proof cites the published fixed-gap input
and supplies the simulation to variable magnitudes. This appendix
remains a paper argument; the separate 713 kernel theorem is a
different construction.

\subsection{Claim and observational contract}
For $\alpha>2$ and $0<a\le 1/2$, let $\mathcal M_{\alpha,a}$ be the probability
laws on $(0,1/2]$ satisfying
\[
 \mu((0,x])\le (x/a)^\alpha\qquad (0<x<a).
\]
A coin has hidden $D\sim\mu$ and $S\mid D=d\sim\operatorname{Ber}(q(d))$;
its successive flips are conditionally independent Bernoulli variables with
parameter $1/2+(2S-1)D$. Fresh coins are iid. All functions $q$ below are measurable,
$0\le q\le1$, and $\theta=\int q\,d\mu\le\rho$, where $0<\rho\le1/4$.
Only actual coin flips are charged. Independent random bits and arithmetic are
free, as in the repository's experiment.

An explicit realization uses mutually independent variables $D_i\sim\mu$,
$U_i\sim\operatorname{Unif}(0,1)$, and $V_{i,n}\sim\operatorname{Unif}(0,1)$,
for $i,n\ge1$, on their countable product probability space. Set
$S_i=\ind_{\{U_i\le q(D_i)\}}$ and
$X_{i,n}=\ind_{\{V_{i,n}\le1/2+(2S_i-1)D_i\}}$.
An additional independent seed supplies every algorithmic random choice.
These are measurable variables and give the required conditional iid flip
tapes and independent fresh coins. The algorithms below use only the next
unused $X_{i,n}$ on a requested coin. Their likelihoods, comparisons, and
stopping decisions depend only on finite observed prefixes. The random-level
and Poisson choices are measurable functions of the independent seed.

Set
\[
\begin{aligned}
 \beta&=1-2/\alpha,\qquad A_0=\frac{\alpha}{\alpha-2},\qquad
 A_1=\frac{\alpha}{(\alpha-2)^2},\\
 w_\alpha&=\frac32\left[
 A_0\left(1+10\log2+\frac1{e\beta}\right)+A_1\right],
 \qquad K=\frac{153}{56}<3.
\end{aligned}
\]
All constants here precede the choice of $a,\mu,\rho,q,B$.

\begin{theorem}[Uniform upper theorem above two]\label{above:thm:upper}
For every $\alpha>2$, $0<a\le1/2$, and $0<\rho\le1/4$, there is one observable
single-coin estimator $Y=Y_{a,\rho}$, independent of $\mu$ and $q$, such that
for every $\mu\in\mathcal M_{\alpha,a}$ and every admissible $q$,
\[
 \E Y=\theta,\qquad \E Y^2\le K\rho,\qquad
 \E T_Y\le w_\alpha/a^2.
\]
In particular all integrability and finite-work assertions are conclusions.
For every $B\ge1$, a family of observable algorithms depending only on
$(\alpha,a,\rho,B)$ has uniform expected work at most $B$ and worst-case RMS
at most
\[
 \min\left\{\frac\rho2,
       \sqrt{K w_\alpha}\frac{\sqrt\rho}{a\sqrt B}\right\}.
\]
The final budgeted estimator need not be unbiased, because clipping or a constant
estimate may be used. The underlying compound-Poisson estimator is unbiased.
\end{theorem}

The theorem gives a known-law result for the exact power law
$F(d)=(d/a)^\alpha$ on $(0,a]$. Its proof also gives the stated uniform upper
bound for the larger margin class. This is not an arbitrary-law solution and
not an estimator independent of a rarity bound. In particular constants diverge
as $\alpha\downarrow2$.

\subsection{Observable asymmetric sequential classifier}
For $k=0,1,\ldots$ define known reference gaps and weights
\[
 r_k=a2^{-(k+1)},\qquad \pi_k=2^{-(k+1)}.
\]
For the observed flips $X_1,X_2,\ldots$ of one coin, let
\[
 L^+_{k,n}=\prod_{i=1}^n(1+2r_k)^{X_i}(1-2r_k)^{1-X_i},\quad
 L^-_{k,n}=\prod_{i=1}^n(1-2r_k)^{X_i}(1+2r_k)^{1-X_i}.
\]
These are likelihood ratios against the fair-coin law. Fix $0<\varepsilon\le1$.
At the first $n\ge1$ for which some $k\ge0$ satisfies either
\[
 L^-_{k,n}\ge\frac1{\pi_k\varepsilon}
 \quad\hbox{or}\quad
 L^+_{k,n}\ge\frac1{\pi_k\rho\varepsilon},                 \tag{B.C}
\]
stop and return $Z_\varepsilon=0$ in the former case and $Z_\varepsilon=1$ in
the latter. Fix negative priority in case both conditions first occur together.
The stopping time is $T_\varepsilon$.

This rule does not observe $D,S$, or $q$. It also does not require infinitely
many operations at any finite observation prefix. If $k\ge n\ge1$, then
$2r_k\le2^{-k}$ and
\[
 L^{\pm}_{k,n}\le(1+2r_k)^n\le\exp(n2^{-k})\le e^{1/2}<2,
 \qquad \frac1{\pi_k\varepsilon}\ge2^{k+1}\ge4.
\]
Thus only the finite set $k<n$ needs to be checked at time $n$. The positive
threshold is even larger. The products for any newly activated index can be
computed from $n$ and the observed head count. Every action and output is
therefore a measurable function of the independent seed and observable history.

\begin{lemma}[Errors and expected work]\label{above:lem:classifier}
For every fixed $0<d\le1/2$ and $s\in\{0,1\}$,
\[
 \Pp(Z_\varepsilon=0\mid d,1)\le\varepsilon,
 \qquad \Pp(Z_\varepsilon=1\mid d,0)\le\rho\varepsilon.
\]
Put $x=\min(d,a)$ and $c_0=1+2\log2$. Then
\[
 \E(T_\varepsilon\mid d,s)
 \le\frac32x^{-2}\left[c_0+\log(a/x)+\log(1/\varepsilon)
                         +s\log(1/\rho)\right].           \tag{B.1}
\]
In particular this stopping time is finite almost surely for every allowed
latent state, including $d=1/2$.
\end{lemma}
\begin{proof}
Under a positive coin, the expected next factor of $L^-_k$ is
$1-4r_kd\le1$; under a negative coin, that of $L^+_k$ is the same.
Consequently these wrong-sign processes are nonnegative supermartingales
starting at one. For completeness, stopping a nonnegative supermartingale
$M$ at its first crossing of $b$ and at a deterministic horizon $N$ gives
$b\Pp(\max_{n\le N}M_n\ge b)\le\E M_{\tau\wedge N}\le1$.
Letting $N\uparrow\infty$ proves the crossing bound $1/b$ without an
unbounded optional-stopping assumption. Applying this to each $k$ and summing
$\sum_k\pi_k=1$ proves the two error bounds, once almost-sure termination is
established below.

Choose $k=\lceil\log_2(a/x)\rceil$. Then $x/4<r_k\le x/2$ (the weaker
$x/4\le r_k\le x/2$ suffices), $r_k\le1/4$, and
$(k+1)\log2\le\log(a/x)+2\log2$.
Under the true sign, the logarithm of the corresponding correct-sign
likelihood ratio is a sum of iid bounded increments, with mean
\[
 g(d,r_k)=\tfrac12\log(1-4r_k^2)
                +d\log\frac{1+2r_k}{1-2r_k}.
\]
The elementary inequalities
$\log(1-z)\ge-z/(1-z)$ and
$\log((1+u)/(1-u))\ge2u$ yield
\[
 g(d,r_k)\ge-\tfrac83r_k^2+4dr_k
            \ge\tfrac83xr_k\ge\tfrac23x^2.               \tag{B.2}
\]
Every upward log increment is at most $\log(1+2r_k)<1$.
If $\sigma$ is the first time this particular log process reaches
\[
 A=\log(1/\varepsilon)+(k+1)\log2+s\log(1/\rho),
\]
then for every integer $N$, its value at $\sigma\wedge N$ is at most $A+1$.
The elementary bounded-horizon identity
\[
 \E\sum_{i=1}^{\sigma\wedge N}\xi_i
   =\sum_{i=1}^N\E[\xi_i\ind_{\{\sigma\ge i\}}]
   =g(d,r_k)\E(\sigma\wedge N)
\]
uses only that $\{\sigma\ge i\}$ is determined before increment $i$.
Hence $\E(\sigma\wedge N)\le(A+1)/g(d,r_k)$.
Monotone convergence proves $\E\sigma\le(A+1)/g(d,r_k)<\infty$ and
therefore $\sigma<\infty$ almost surely. Since $T_\varepsilon\le\sigma$,
equations (B.1)--(B.2) follow. This also completes the termination needed for
the error assertions.
\end{proof}

\subsection{The moment calculation that removes the logarithmic loss}
The relevant distinction from symmetric high-confidence screening is that
$\log(1/\rho)$ in (B.1) is paid only on positive coins.

\begin{lemma}[Margin integration]\label{above:lem:integration}
If $\mu\in\mathcal M_{\alpha,a}$, $x=\min(D,a)$, $\alpha>2$, and
$0\le q\le1$, $\E q(D)\le\rho$, then
\[
\begin{aligned}
 \E x^{-2}&\le A_0a^{-2},\\
 \E[x^{-2}\log(a/x)]&\le A_1a^{-2},\\
 \E[q(D)x^{-2}]&\le A_0a^{-2}\rho^\beta.
\end{aligned}                                           \tag{B.3}
\]
For the exact power law the first two inequalities are equalities, and the
last is attained by $q(d)=\ind_{\{d\le a\rho^{1/\alpha}\}}$.
\end{lemma}
\begin{proof}
The variable $x$ stochastically dominates the power-law variable $aU^{1/\alpha}$,
where $U$ is uniform on $(0,1)$. Integrating decreasing nonnegative functions,
or equivalently applying the layer-cake identity to their sublevel events,
gives the first two inequalities. Their explicit integrals are
\[
 \frac\alpha{a^\alpha}\int_0^a d^{\alpha-3}\,dd=A_0a^{-2},
 \qquad
 \frac\alpha{a^\alpha}\int_0^a d^{\alpha-3}\log(a/d)\,dd=A_1a^{-2}.
\]
For the third, Tonelli's theorem and $q\le1$ give
\[
 \E[qx^{-2}]
 =a^{-2}\E q+\int_{a^{-2}}^\infty\E[q\ind_{\{x^{-2}>u\}}]\,du
 \le a^{-2}\rho+\int_{a^{-2}}^\infty
                  \min\{\rho,(a^2u)^{-\alpha/2}\}\,du.
\]
Splitting at $a^{-2}\rho^{-2/\alpha}$ evaluates the last expression as
$A_0a^{-2}\rho^{1-2/\alpha}$. The claimed equality case follows by direct
integration. These arguments concern nonnegative extended integrals;
finiteness is established by the displayed evaluations.
\end{proof}

Combining (B.1) and (B.3), with $\theta\le\rho$, gives the uniform bound
\[
 \E T_\varepsilon\le\frac{3}{2a^2}
 \left[A_0\{c_0+\log(1/\varepsilon)\}+A_1
                  +A_0\rho^\beta\log(1/\rho)\right].     \tag{B.4}
\]
Here $\rho^\beta\log(1/\rho)\le1/(e\beta)$, obtained by maximizing
$te^{-\beta t}$ over $t\ge0$. This is where $\alpha>2$ is essential.
The logarithm is controlled uniformly in the supplied rarity bound.

\subsection{Unbiased random-level estimator}
The single-term randomized telescoping device is an established unbiased
multilevel ingredient; see Vihola~\cite[Theorem 3]{vihola}. The particular classifier and
the rarity-sensitive moment and work calculation below are analyzed explicitly.
Let $\varepsilon_j=16^{-(j+1)}$ for $j\ge0$ and use the above classifiers
$Z_j=Z_{\varepsilon_j}$ on one common coin-flip tape. Define
\[
 \Delta_0=Z_0,\qquad \Delta_j=Z_j-Z_{j-1}\quad(j\ge1).
\]
Choose an independent $J$ with $p_j=\Pp(J=j)=2^{-(j+1)}$, draw one fresh coin,
and return $Y=\Delta_J/p_J$. Only the classifiers actually needed for this
increment are evaluated. Both threshold families in (B.C) grow by the same
factor $16$ between successive levels, so $T_j\ge T_{j-1}$ pathwise. The
increment costs exactly $T_J$, with the earlier decision stored from the same
prefix. No assumption about independence between classifier levels is made.

The conditional error bounds give
\[
 \Pp(Z_j\ne S)\le (\theta+\rho)\varepsilon_j\le2\rho\varepsilon_j,
 \qquad \E Z_0^2=\E Z_0\le\theta+\rho\varepsilon_0\le\tfrac{17}{16}\rho.
\]
For $j\ge1$, a disagreement implies an error of at least one classifier, hence
\[
 \E\Delta_j^2\le2\rho(\varepsilon_j+\varepsilon_{j-1})
                    =\tfrac{17}{8}\rho\,16^{-j}.
\]
All square terms are nonnegative, so Tonelli gives
\[
 \E Y^2=\sum_{j\ge0}\frac{\E\Delta_j^2}{p_j}
 \le\left(\frac{17}{8}+\frac{17}{4}\sum_{j\ge1}8^{-j}\right)\rho
 =\frac{153}{56}\rho.                                  \tag{B.5}
\]
This proves square integrability, hence integrability, before using signed
expectations. Alternatively $\sum_j\E|\Delta_j|<\infty$ directly follows
from the same disagreement bounds because $|\Delta_j|=\Delta_j^2$ for $j\ge1$.
Consequently
\[
 \E Y=\sum_{j\ge0}\E\Delta_j
       =\lim_{j\to\infty}\E Z_j=\E S=\theta.             \tag{B.6}
\]
Finally (B.4), $\E(J+1)=2$, and $\log(1/\varepsilon_j)=4(j+1)\log2$ imply
\[
 \E T_Y\le\frac3{2a^2}
 \left[A_0\{1+10\log2+1/(e\beta)\}+A_1\right]
 =w_\alpha/a^2.                                         \tag{B.7}
\]
Equations (B.5)--(B.7) prove the single-coin part of Theorem~\ref{above:thm:upper}, with
separate uniform second-moment and work bounds.

\subsection{Exact expected-budget calibration and matching power-law rate}
Independent Poisson replication and its variance formula are also established
ingredients; see Vihola~\cite[Example 8]{vihola}. We give the direct calculation
to verify the exact expected budget, including when the expected count is below one.
Set $\lambda=Ba^2/w_\alpha>0$. Draw an independent
$M\sim\operatorname{Poisson}(\lambda)$ and run independent copies $Y_1,\ldots,Y_M$
of the single-coin procedure on fresh coins. Return
\[
 \widehat\theta=\lambda^{-1}\sum_{i=1}^M Y_i.
\]
The empty sum is zero. Conditioning first on $M$, then using its mean and
variance $\lambda$, proves
\[
 \E\widehat\theta=\theta,
 \quad \Var(\widehat\theta)=\lambda^{-1}\E Y^2
            \le\frac{K w_\alpha\rho}{a^2B},
 \quad \E T\le\lambda w_\alpha/a^2=B.
\]
All sums stop almost surely and have finite expected work. These identities are
valid even when $\lambda<1$; no rounding of the budget is hidden. Clip the result
to $[0,\rho]$, which cannot increase squared error for $\theta\in[0,\rho]$.
When its displayed RMS bound exceeds $\rho/2$, instead output $\rho/2$ at zero
work. The choice uses only the allowed parameters. This proves the budgeted
part of Theorem~\ref{above:thm:upper}.

\begin{lemma}[Localized fixed-gap RMS lower bound]\label{above:lem:fixedlower}
There is a universal $c>0$ such that the minimax RMS risk for fixed gap
$0<a\le1/2$, prior $0\le\vartheta\le\rho\le1/4$, and uniform expected
work at most $B\ge1$ is at least
$c\min\{\rho,\sqrt\rho/(a\sqrt B)\}$.
This lemma uses Lee--Valiant Corollary 19 and Lemma 18 as external inputs.
\end{lemma}
\begin{proof}
The exact imported assertions imply the following: there are universal
$a_*,c_*,C_*>0$ such that, if $a<a_*$ and $u,v,v/u<c_*$, the KL divergence
of any finite-prefix fully adaptive fixed-gap transcript under sign priors
$u$ and $u+v$ is at most $C_*v^2a^2\E_u T/u$.
Indeed, Corollary 19 bounds each single-coin hybrid's squared Hellinger
distance by a universal multiple of $v^2a^2/u$ times its expected observations
under $u$. Lemma 18 bounds joint KL by the sum of those distances. Summing
coin work counts each observation once. Their Hellinger normalization is
absorbed in $C_*$. We use the original-prior, not midpoint-prior, expectation.
All their smallness hypotheses are enforced in the choices below.

For arbitrary almost-surely terminating policies with finite expected work,
first stop at $T\wedge n$ and number coins by first observation. At most $n$
coins are then observed, so the finite-prefix statement applies. This numbering
preserves the paper's creation and interleaving model: an unobserved coin has
provided no information, and its iid latent draw can be deferred until its first
observation. A dictionary from the policy's original coin labels to the new
indices preserves every revisit. Unobserved-label choices use only the seed
and observed history, so this relabeling changes neither their law nor any
charged observation or output. Include the
independent random seed in the transcript; conditioning on it and averaging
extends the deterministic statement. Relative entropy along increasing prefix
sigma fields converges to full transcript relative entropy. The stopped output
is measurable in this sigma field. Thus the same bound holds at unbounded
stopping, without a globally fixed finite coin family.

Write $r=\min\{\rho,\sqrt\rho/(a\sqrt B)\}$, and take
$u=b\rho$, $v=\gamma r$ for fixed sufficiently small positive universal
constants $b,\gamma$. Choose $b+\gamma\le1$ and $\gamma/b<c_*$,
so $u,u+v\in[0,\rho]$ and all prior smallness assumptions hold.
Since $r^2a^2B\le\rho$, for $a<a_*$ the transcript KL is at most
$C_*\gamma^2/b$.

For $a\ge a_*$, give the algorithm each coin's true sign on its first
observation. This more informative experiment simulates all requested flips
from that sign. Each newly observed sign has Bernoulli KL at most $2v^2/u$,
using $u+v\le1/4$ and the KL--chi-square inequality. The expected number of
newly observed signs is at most $\E_u T\le B$, so full KL is at most
$2v^2B/u\le2\gamma^2/(ba_*^2)$. The same finite-prefix argument justifies
unbounded stopping. This also covers the endpoint $a=1/2$.

Choose $\gamma$ small enough that both displayed KL bounds are at most $1/8$.
Pinsker's inequality then bounds total variation by $1/4$. If an estimator has
RMS at most $R$ at both priors, thresholding halfway between them gives a test
whose two error probabilities sum to at most $8R^2/v^2$. Every test has error
sum at least $1-\mathrm{TV}\ge3/4$. Therefore $R\ge v/4$, proving the claim
with $c=\gamma/4$. This proves all budget and rarity regimes, including the
saturated regime, from the precise source inputs rather than an informal
high-confidence asymptotic statement.
\end{proof}

\begin{corollary}[Exact-power regime $\alpha>2$]\label{above:cor:power}
Let $\mu$ have distribution function $(d/a)^\alpha$ on $(0,a]$ with $\alpha>2$.
Subject to the published fixed-gap lower bound of Lee--Valiant, there is a
universal $c>0$ and, for each $\alpha>2$, a finite $C_\alpha$ such that for all
$0<a\le1/2$, $0<\rho\le1/4$, and $B\ge1$,
\[
 c\min\left\{\rho,\frac{\sqrt\rho}{a\sqrt B}\right\}
 \le R_{\mu,\rho}(B)
 \le C_\alpha\min\left\{\rho,\frac{\sqrt\rho}{a\sqrt B}\right\}.
\]
Moreover this risk is comparable, with constants depending only on $\alpha$,
to $\Psi_{\mu,\rho}(B^{-1/2})$.
\end{corollary}
\begin{proof}
The upper bound was just proved. For the lower bound it is enough to restrict
to constant sign functions $q(d)=\vartheta\in[0,\rho]$. An experiment using
fixed-gap $a$ coins can simulate the variable-gap experiment at the same charged
work: sample $D\sim\mu$ privately for each fresh fixed-gap coin; on each request,
flip that coin once, and with probability $D/a$ return its flip, otherwise
return an independent fair bit. Conditional on the fixed coin's sign and its
privately sampled $D$, the
returned flips are iid Bernoulli$(1/2\pm D)$. The same latent $D$ is retained
on revisits, and is not disclosed to the simulated algorithm. The simulation
uses one charged flip per returned observation even when that flip is discarded.
It preserves all interleaving and stopping choices. Thus a smaller variable-gap
risk would contradict Lemma~\ref{above:lem:fixedlower}. This reduction does not need
the continuation note's arbitrary-law lower-bound theorem.

For comparison with the functional, constant allocation $w\equiv\rho$ and
$D\le a$ give
\[
 \Psi_{\mu,\rho}(h)\ge\min\{\rho,h\sqrt\rho/a\}.
\]
For every feasible $w$, Cauchy--Schwarz and $\E D^{-2}=A_0/a^2$ give
\[
 \int\min\{w,h\sqrt w/D\}\,d\mu
 \le\min\{\rho,h\sqrt{\rho A_0}/a\}.
\]
Taking the supremum proves the claimed comparison, including saturation.
\end{proof}

\subsection{Scope, edge cases, and failed approaches}
The classifier and estimator operate at $D=1/2$, with no likelihood logarithm
at zero because the reference gaps satisfy $r_k\le1/4$. At arbitrarily small
$D>0$ each conditional expected stopping time is finite; the margin assumption
then proves population integrability. The level $J$ has unbounded support and
finite mean. For the admissible law $\mu=\delta_{1/2}$, $a=1/2$, and $q=0$,
all observations are zero. At time $n$ every correct-sign likelihood is at most
$(3/2)^n$, whereas every level-$j$ correct-sign threshold is at least
$2/\varepsilon_j$. Thus each fixed $n$ is exceeded for sufficiently large $j$.
Since all $p_j>0$, actual work has genuinely unbounded support and finite mean,
not merely an unused unbounded seed. A finite Poisson number of such procedures
also has finite expected work.

The global constraint $\int q\,d\mu\le\rho$ is used precisely in (B.3); no
pointwise bound $q\le\rho$ is imposed. Its extremizer puts all positive mass
on the smallest gaps. The algorithm never assumes known positive or negative
conditional magnitude laws. The result is uniform jointly in $B,\rho,a$ for
each fixed $\alpha>2$, but is not uniform as $\alpha\downarrow2$.

Two earlier approaches exposed the obstruction addressed here. Symmetric
high-confidence screening charges $\log(1/\rho)$ on all coins, retaining an
unwanted logarithm. An unamplified ballot-survival estimator has positive
conditional second moment of order $1/D$; placing positives near the threshold
then gives an unwanted $\rho^{1-1/\alpha}$ moment scale. The asymmetric
classifier followed by an unbiased random-level correction avoids both
problems. The negative classifier error is $\rho\varepsilon$, whereas the
positive error is only $\varepsilon$; the extra positive work is absorbed by
the sharp global-rarity integral in (B.3).

The same proof does not settle $\alpha=2$: its inverse-square integral diverges,
and the uniform control of $\rho^\beta\log(1/\rho)$ loses meaning at $\beta=0$.
It does not characterize arbitrary laws with a finite inverse-square moment,
because that moment alone does not uniformly control the logarithmically
weighted moment or the rarity-weighted cost. It does not refute that larger
conjecture. A direct finite-support modification gives finite bounds with
support-dependent constants, but this is not the universal two-atom theorem.
